% index.tex
% -----------------------------------------------------------------------------
% Master index for the atomic-unit decomposition worked examples.
%
% Each example lives in its own self-contained, independently-compilable file
% (example-1-damages.tex ... example-5-renda.tex). This index links to the
% compiled PDF of each example. Build the examples first, then this index:
%
%   for f in example-1-damages example-2-biography example-3-narrative \
%            example-4-summary example-5-renda index; do
%     pdflatex -interaction=nonstopmode "$f.tex"; pdflatex -interaction=nonstopmode "$f.tex"
%   done
%
% The links open the sibling PDFs, so keep index.pdf next to the example PDFs.
% -----------------------------------------------------------------------------
\documentclass[11pt]{article}

\usepackage[margin=1in]{geometry}
\usepackage{enumitem}
\usepackage{xcolor}
\usepackage[colorlinks=true,urlcolor=blue!60!black,linkcolor=blue!60!black]{hyperref}

\title{Atomic-Unit Decomposition and Relation Inference\\\large Worked Examples --- Master Index}
\author{FactReasoner --- Ideation}
\date{\today}

% Convenience macro: \example{file}{Title}{Description}
\newcommand{\examplelink}[3]{%
  \item \textbf{\href{run:./#1.pdf}{#2}}\\%
  \nolinkurl{#1.tex}\, $\rightarrow$ \,\nolinkurl{#1.pdf}\\[2pt]%
  #3\par\medskip%
}

\begin{document}
\maketitle

\section*{Overview}

This collection demonstrates a single recurring task applied to several kinds of
text: decompose a paragraph into \emph{self-contained atomic units}, then infer
two relation types between them ---

\begin{itemize}[leftmargin=*,nosep]
  \item \textbf{logical prerequisite} ($A \rightarrow B$): $A$ must be true for
        $B$ to be true or meaningful;
  \item \textbf{position invalidation / contradiction} ($A \neq B$): $A$, true at
        its position in the text, contradicts a claim $B$ appearing later (or a
        later atom overrides an earlier one).
\end{itemize}

Each example is a standalone document with its source paragraph, atomic units,
relation lists, and a TikZ relation graph. Click a title to open that example's
PDF (build the examples first; see the header of \nolinkurl{index.tex}).

\section*{Examples}

\begin{enumerate}[leftmargin=*]

  \examplelink{example-1-damages}{1.\ A legal damages paragraph}{%
    A contract-damages paragraph in which a mid-text ruling that the agreement
    was \emph{void ab initio} retroactively invalidates the earlier
    breach-of-contract claims, while a separate unjust-enrichment award survives.
    Introduces both relation types and the pivotal-invalidation pattern.}

  \examplelink{example-2-biography}{2.\ A biography (with planted contradictions)}{%
    The Lanny Flaherty biography: first as an internally consistent passage
    (pure prerequisite / subsumption structure, no contradictions), then a
    modified version with several planted contradictions (birth year, birthplace,
    nationality, career start, voice) drawn as pairwise invalidation edges.}

  \examplelink{example-3-narrative}{3.\ A narrative passage (Elinor)}{%
    A scene from \emph{Sense and Sensibility} with a late revelation
    (Robert, not Edward, married Lucy) that invalidates the passage's opening
    premise. Includes a ground-truth re-decomposition, an alignment table, and a
    temporal-order discrepancy plot comparing narration order to chronology.}

  \examplelink{example-4-summary}{4.\ Synthesizing a summary from a reliable and an unreliable source}{%
    A summary woven from an accurate source A and a sensational, error-laden
    source B (the Keyara Jones example), engineered so later reliable atoms
    positionally invalidate earlier false ones, including a self-contradictory
    statistic.}

  \examplelink{example-5-renda}{5.\ A real legal case (\emph{R v Renda})}{%
    The lead appeal in \emph{Renda, R v} [2005] EWCA Crim 2826. Two summaries of
    the same appeal: an engineered ordering (S) that manufactures positional
    contradictions, and a faithful, naturally-ordered summary (K) whose relations
    collapse to a pure implication chain --- illustrating that positional
    invalidation depends on how a text is arranged, not only on its facts.}

\end{enumerate}

\section*{Relation-graph conventions}

Across the examples the graphs share a common visual language: solid blue arrows
are logical prerequisites ($A \rightarrow B$); dashed red arrows are
contradictions / position invalidations ($A \neq B$); pivotal or
holding/revelation nodes are highlighted (orange), false or invalidated atoms are
tinted red, and reliable or surviving atoms are tinted green.

\end{document}
